% Requires \usepackage{booktabs}. Full-width floats for two-column conference layouts.
% Source tables: T1
\begin{table*}[t]
\centering
\small
\setlength{\tabcolsep}{4pt}
\caption{Comparison on the project VisDrone test split under a common evaluator.}
\label{tab:selected-1}
\begin{tabular}{lrrrrrrr}
\toprule
Method & Params (M) & AP & AP50 & AP75 & APs & APm & APl \\
\midrule
D-FINE-L (retrained) & 31.146 & 22.679 & 38.940 & 22.713 & 13.478 & 32.226 & 40.858 \\
Ours (Full) & 35.244 & 24.229 & 42.009 & 24.208 & 14.618 & 34.290 & 44.873 \\
Delta (pp) & -- & 1.550 & 3.069 & 1.495 & 1.140 & 2.064 & 4.016 \\
\bottomrule
\end{tabular}
\par\smallskip
\begin{minipage}{0.97\textwidth}\footnotesize AP values are percentages. Both checkpoints were selected by native validation AP and evaluated on the same 1,610 project test images using YOLO-derived COCO annotations, 640 input, FP32 and maxDet=300, without TTA or extra NMS. Training recipes differ: D-FINE uses batch 4, AdamW and FP32; ours uses batch 8, MuSGD and AMP. Parameters are registered model totals, not deployment-pruned counts. This is not an official challenge ranking.\end{minipage}
\par\medskip
\end{table*}

% Source tables: T2
\begin{table*}[t]
\centering
\small
\setlength{\tabcolsep}{4pt}
\caption{Component comparison on native VisDrone validation results.}
\label{tab:selected-2}
\begin{tabular}{lrrrrrrr}
\toprule
Model & G & B & F & P & R & AP50 & AP \\
\midrule
Base & -- & -- & -- & 46.761 & 41.731 & 38.663 & 21.911 \\
G & + & -- & -- & 59.480 & 49.366 & 49.382 & 29.620 \\
G+F & + & -- & + & 59.391 & 50.049 & 49.548 & 29.773 \\
G+B & + & + & -- & 59.826 & 48.524 & 48.783 & 29.277 \\
Full & + & + & + & 60.482 & 49.693 & 49.890 & 30.109 \\
\bottomrule
\end{tabular}
\par\smallskip
\begin{minipage}{0.97\textwidth}\footnotesize G: LRS-GFDR; B: Capacity-BPDD-v3-Safe; F: P3-only FIA. All metrics in a row are from the same validation-AP-best epoch. Base, G and G+F are registered reference runs; G+B and Full share recorded data, initialization and training settings. Source-path equivalence remains to be audited. The table is not a fully verified factorial experiment.\end{minipage}
\par\medskip
\end{table*}

% Source tables: T3, T4
\begin{table*}[t]
\centering
\small
\setlength{\tabcolsep}{4pt}
\caption{Internal design comparisons on native VisDrone validation results.}
\label{tab:selected-3}
\textbf{(a) Localization design}\par\smallskip
\begin{tabular}{lrrrr}
\toprule
Variant & Geometry & LRS & AP50 & AP \\
\midrule
Original FDR & reference & reference & 48.473 & 29.078 \\
Clean-FDR & -- & -- & 49.320 & 29.696 \\
Clean-FDR + geometry & + & -- & 49.130 & 29.361 \\
LRS-GFDR & + & + & 49.382 & 29.620 \\
\bottomrule
\end{tabular}
\par\smallskip
\begin{minipage}{0.97\textwidth}\footnotesize These are registered design variants, not a fully source-matched toggle experiment. Geometry reduces validation AP relative to Clean-FDR in these runs; LRS recovers part of that difference. The table does not establish a standalone accuracy gain from geometry or prove a second-dataset stability claim. Reference denotes switches not mapped to the current implementation.\end{minipage}
\par\medskip
\textbf{(b) BPDD recipe and FIA compatibility}\par\smallskip
\begin{tabular}{lrrr}
\toprule
BPDD recipe & FIA & AP50 & Best AP \\
\midrule
Capacity v2 & -- & 49.620 & 29.848 \\
Capacity v2 & + & 49.585 & 29.736 \\
Capacity v3-Safe & -- & 48.783 & 29.277 \\
Capacity v3-Safe & + & 49.890 & 30.109 \\
\bottomrule
\end{tabular}
\par\smallskip
\begin{minipage}{0.97\textwidth}\footnotesize The v2-to-v3-Safe change is a recipe bundle, not an isolated toggle for each distillation or scheduling choice. The v3-Safe BPDD-only run is worse than v2, while its joint run is better. These results motivate compatibility-oriented design rather than universally improved standalone BPDD.\end{minipage}
\par\medskip
\end{table*}
